\section{Discussion}
\label{sec:discussion}

The results of \cref{sec:results} describe a policy that tracks task-space poses to a few centimetres under lunar gravity while keeping a user-specified constraint set largely satisfied. This section examines the mechanisms behind that behaviour, what the constrained formulation contributes to it, and the boundaries of what the study establishes.

\subsection{What the Constraint Set Contributes}

The ablation of \cref{exp:constraints} shows that the observed motion style comes from the constraint set, not from the task reward. Removing all soft constraints leaves tracking performance broadly intact while allowing knee torque violations to rise from $0.006\%$ to $13.92\%$ of episode time and joint deviation violations from $6.45\%$ to $38.60\%$. The policy, in other words, can solve the tracking task perfectly well by operating outside the intended envelope; nothing in the reward prevents it. What the constraints supply is not task competence but admissibility, and they supply it in physical units, so nothing needs retuning when the task changes.

The per-constraint ablations sharpen this. Removing $c_{dev}$ raises its own violation time eightfold while leaving the torque and foot-contact constraints essentially unchanged, and removing $c_v$ roughly doubles base velocity violations without disturbing the rest. Constraints act with a degree of independence, which is a practical property: the set can be extended or relaxed to reflect a different platform or mission profile without the compensating retuning that reward-shaped safety terms typically demand.

The residual violation rates deserve interpretation. The base velocity constraint is violated for $11.58\%$ of the episode by the baseline policy, the highest rate in the set, and this follows directly from the structure of the problem. The $0.25\,\mathrm{m/s}$ limit and the $10\,$s episode together permit exactly the $2.5\,$m maximum commanded displacement, so the most distant commands are reachable only by translating at the limit with no margin. A policy that never violated $c_v$ would necessarily forfeit those commands. Here the formulation arbitrates a real conflict between task completion and motion style, and the bounded termination probability of a soft constraint is the mechanism that lets it resolve in favour of the task when the alternative is failure. This is the intended behaviour of the relaxation described in \cref{sec:cat}, not a defect in its realisation.

\subsection{Reward Design under a Multiplicative Objective}

Three elements of the reward formulation proved consequential, and the reasons generalise beyond this application.

The first is the requirement that every reward term be non-negative. \cref{sec:positive-rewards} establishes why: the modulation $(1-\delta)$ attenuates penalties and task rewards alike, so a penalty-form regularizer becomes cheaper exactly when a constraint is violated. This is not a peripheral concern, because the quantities one naturally wishes to regularize, joint power and action rate, are dynamically coupled to the quantities the constraints bound. Reformulating them as bounded positive rewards preserves the policy ordering they induce under modulation. Any application of termination-based constraint enforcement faces this issue, and the design rule transfers directly.

The second is the implicit phase switch. The multiplicative composition $r_{\text{pose}} \cdot r_{\text{base}}$ with a saturating proximity term produces the locomotion-then-manipulation sequence without a state machine, a gate, or a scheduled transition: far from the target essentially all gradient flows through $r_{\text{base}}$, and inside the proximity region that term saturates and the gradient passes entirely to $r_{\text{pose}}$. Saturation is what makes this work. A monotonically increasing proximity term would not: clipping renders base placement inside the region reward-neutral and leaves the policy free to position the base wherever the arm is best conditioned.

The third is the error-adaptive kernel. Commanded poses span three orders of magnitude of position error over an episode, and no fixed kernel width serves that range. A wide kernel cannot discriminate a $2\,$cm error from a $10\,$cm one, and a narrow kernel is numerically flat at the distances where the policy must decide where to walk. Contracting the kernel with the error keeps the reward in a well-conditioned band throughout, and does so without scheduling on training iteration, so the sharpening follows the policy's actual competence instead of a fixed timetable.

\subsection{Behaviour under Lunar Conditions}

The stability formulation is where reduced gravity most visibly shapes the method. With $\lVert\boldsymbol{g}\rVert \approx 1.63\,\mathrm{m/s^2}$, the accelerations produced by swing-leg and manipulator motion are a substantial fraction of the gravitational term, so the net acceleration the body must resist through its contacts departs materially from the vertical during ordinary motion. A static support-polygon criterion is blind to this, whereas the gravito-inertial margin of \cref{sec:stability} measures the quantity that actually governs tipping. Using the mean over polyhedron faces rather than the minimum converts a non-smooth metric into a dense signal that rewards margin in all directions, and the affine mapping onto the unit interval keeps it compatible with the positivity requirement above.

The robustness results indicate that terrain geometry, not substrate, dominates the error budget. Tracking error is flat across regolith densities, rising by roughly $2\,$cm only where loose soil coincides with steep hills, where reduced friction and slope compound. Across terrain morphology the picture differs: rough terrain leaves accuracy unaffected, hills degrade it mildly, and craters at $30^\circ$ raise mean error to $13\,$cm. That the degradation appears only at the extreme of the difficulty range, and does not build progressively across it, suggests the policy holds one strategy until slope approaches the platform's traversability limit. Since the compliant contact model was withheld from training entirely, the flatness across soil types is also evidence that the policy tolerates this class of model mismatch instead of having memorised one contact regime.

\subsection{Scope and Limitations}

Several boundaries of the study follow from choices made in framing it.

The evaluation is conducted entirely in simulation, and the results characterise the framework under the environment model of \cref{sec:lunar-design} rather than on hardware. Contact with granular regolith is represented by a compliant linear model, which reproduces energy dissipation and softened contact transitions but not the constitutive behaviour of a granular medium; foot-soil interaction under sustained loading would require a terramechanics or discrete-element treatment. Terrains are procedurally generated, so the generalisation demonstrated is within that family and does not extend to geometries drawn from measured lunar topography. Perception is modelled as an idealised elevation measurement with Gaussian noise, which does not represent the structured failure modes (e.g shadowing, saturation and dust contamination) expected at high lunar latitudes.

Our comparisons are internal. The constraint ablations isolate the contribution of the constraint set, and the terrain and soil sweeps characterise sensitivity, but we do not evaluate decoupled or unconstrained architectures under identical conditions; the case for unified whole-body control rests here on the task formulation rather than on a controlled measurement.

Two extensions follow naturally. The task is free-space pose tracking, and contact-rich operations such as coring or fastening would require wrist force sensing, an interaction-force objective, and extension of the constraint set to the arm, which is presently unconstrained on the grounds that it operates out of contact. Separately, the piecewise-constant kernel schedule introduces discontinuities at its interval boundaries; a continuously varying $\sigma(e)$ would retain the conditioning benefit while restoring monotonicity of the reward in the tracking error.
